\chapter{Thrush}
\index{Thrush}

\medskip
\noindent\textit{Educational reference; always seek professional advice for individual care.}

\section*{Definition and Overview}
Thrush is a common name for candidiasis, an overgrowth of Candida yeasts---most often Candida albicans---on mucous membranes or moist skin. In everyday local usage, thrush often means vaginal candidiasis, but the same organism causes oral thrush (white patches in the mouth), nipple thrush in breastfeeding, male genital candidiasis, nappy-area infection, and, in vulnerable people, oesophageal or invasive disease. Candida commonly lives harmlessly on the body; symptoms appear when the local environment or immunity shifts. Most mucosal thrush is uncomfortable rather than dangerous and responds to topical or short-course antifungal treatment, but recurrent disease needs a fuller assessment. This chapter covers typical presentations, self-care, pharmacy and GP pathways, and situations requiring urgent care.

\section*{Epidemiology}
Vaginal thrush affects a large proportion of women at least once; a smaller group experiences recurrent vulvovaginal candidiasis (four or more symptomatic episodes yearly). Oral thrush is seen in infants, denture wearers, people using inhaled corticosteroids, those with dry mouth, diabetes, or HIV, and after broad-spectrum antibiotics. Skin-fold candidiasis increases with obesity and hot humid weather. Invasive candidiasis is mainly a hospital problem linked to central lines, ICU care, neutropenia, and major abdominal surgery. Antifungal resistance, including non-albicans species, is an emerging stewardship concern when thrush is treated repeatedly without diagnosis.

\section*{Aetiology and Pathophysiology}
\begin{itemize}
    \item \textbf{Yeast overgrowth:} Candida switches from commensal yeast to invasive hyphal forms that irritate epithelium
    \item \textbf{Antibiotic effect:} loss of protective bacterial flora after antibiotics is a classic trigger for vaginal and oral thrush
    \item \textbf{Hormonal and pregnancy factors:} higher glycogen and immune shifts favour vaginal candidiasis
    \item \textbf{Diabetes and hyperglycaemia:} glucose-rich secretions support yeast growth
    \item \textbf{Immunosuppression and inhaled steroids:} reduce local or systemic control of Candida
    \item \textbf{Moisture and occlusion:} nappies, tight synthetic clothing, and skin folds create a niche
    \item \textbf{Sexual activity:} can irritate and transfer organisms but thrush is not classified as a classic STI
    \item \textbf{Dentures and xerostomia:} continuous carriage under plates
\end{itemize}

\section*{Clinical Features}
\begin{itemize}
    \item \textbf{Vaginal thrush:} itch, burning, dysuria externally, thick white ``cottage cheese'' discharge, non-offensive compared with bacterial vaginosis; dyspareunia
    \item \textbf{Oral thrush:} creamy white plaques on tongue or cheeks that may scrape off leaving red base; taste change; angular cheilitis at mouth corners
    \item \textbf{Male genital thrush:} red itchy penile glans, white patches, discomfort after sex
    \item \textbf{Skin-fold candidiasis:} beefy red rash with satellite pustules under breasts, in groin, or abdominal folds
    \item \textbf{Nappy candidiasis:} persistent bright rash with satellites not responding to simple barrier cream
    \item \textbf{Oesophageal candidiasis:} painful swallowing and retrosternal pain in immunocompromised people
    \item \textbf{Nipple/ductal symptoms in breastfeeding:} sharp feeding pain out of proportion to visible trauma (diagnosis can be tricky)
\end{itemize}

\section*{Differential Diagnosis}
\begin{itemize}
    \item \textbf{Bacterial vaginosis and trichomoniasis:} different discharge and odour patterns; need proper STI-aware assessment when uncertain
    \item \textbf{Genital herpes and contact dermatitis:} ulcers or allergen history
    \item \textbf{Lichen sclerosus and other dermatoses:} chronic white plaques, scarring---not yeast alone
    \item \textbf{Oral leukoplakia and lichen planus:} white patches that do not scrape like thrush
    \item \textbf{Geographic tongue and milk residue in babies:} innocuous mimics
    \item \textbf{Urinary tract infection:} internal dysuria and frequency more than external itch
    \item \textbf{Cytology of ``recurrent thrush'':} many women are misattributed; examine and test before endless antifungals
\end{itemize}

\section*{When to Seek Medical Care}
First typical episode in adults can often be managed with pharmacist advice. See a GP if you are pregnant, symptoms are atypical or recurrent, over-the-counter treatment fails, you have diabetes or immunosuppression, there is oral thrush in adults without clear cause, or a child is feeding poorly.

\textbf{Contact your local emergency services for:}
\begin{itemize}
    \item Severe systemic illness, high fever, and widespread rash in an immunocompromised person
    \item Difficulty breathing or swallowing with severe oral or throat swelling
    \item Suspected invasive infection in hospitalised or neutropenic patients with shock
\end{itemize}

\section*{Diagnosis and Investigations}
\begin{itemize}
    \item \textbf{Clinical recognition:} often sufficient for classic vaginal or infant oral thrush
    \item \textbf{Vaginal pH and microscopy/culture:} when recurrent, resistant, or diagnosis unsure
    \item \textbf{STI testing:} if risk factors or mixed symptoms
    \item \textbf{Blood glucose / HbA1c:} in recurrent disease
    \item \textbf{HIV and immune work-up:} unexplained adult oral or oesophageal candidiasis
    \item \textbf{Oral swab or endoscopic confirmation:} for oesophageal disease
    \item \textbf{Species identification and susceptibility:} selected refractory cases
\end{itemize}

\section*{Management}
\begin{itemize}
    \item \textbf{Topical azoles:} clotrimazole or miconazole creams and pessaries for vaginal disease; courses as labelled
    \item \textbf{Single-dose oral fluconazole:} common for uncomplicated vaginal thrush in non-pregnant adults when appropriate
    \item \textbf{Pregnancy:} topical treatments preferred; avoid fluconazole unless specialist advice says otherwise
    \item \textbf{Oral nystatin or miconazole gel:} infant and adult oral thrush; sterilise teats and treat mother if breastfeeding co-infection suspected
    \item \textbf{Inhaler technique:} spacer and mouth rinse after corticosteroid inhalers
    \item \textbf{Skin-fold care:} keep dry, topical antifungal, address weight and clothing
    \item \textbf{Recurrent vaginal candidiasis:} induction then maintenance regimens under GP or gynaecology guidance; partner treatment only if male symptoms
    \item \textbf{Systemic antifungals for invasive disease:} hospital protocols, line removal, ophthalmology checks as indicated
\end{itemize}

\section*{Complications}
\begin{itemize}
    \item \textbf{Chronic itch--scratch cycle and secondary bacterial infection}
    \item \textbf{Psychological and sexual distress:} especially with recurrent genital thrush
    \item \textbf{Cracked nipples and early weaning:} if breastfeeding pain is unaddressed
    \item \textbf{Oesophageal and invasive candidiasis:} morbidity in immunocompromised hosts
    \item \textbf{Azole resistance and non-albicans species:} after repeated empiric therapy
    \item \textbf{Missed alternative diagnoses:} lichen sclerosus, STI, or pre-cancerous lesions treated as ``thrush'' for months
\end{itemize}

\section*{Prognosis and Outcomes}
Uncomplicated thrush usually clears quickly with correct treatment. Recurrence is the main challenge and often improves when triggers (antibiotics, uncontrolled glucose, steroids without rinsing) are managed. Invasive disease outcomes depend on host immunity and speed of antifungal therapy. Most healthy infants with oral thrush do well.

\section*{Prevention and Self-Care}
\begin{itemize}
    \item Avoid unnecessary antibiotics; complete necessary courses but report thrush early
    \item Rinse mouth after steroid inhalers; clean dentures nightly and leave them out while sleeping when possible
    \item Choose breathable cotton underwear; change out of wet gym clothes promptly
    \item Wipe front to back; avoid douching and heavily fragranced genital products
    \item Optimise diabetes control
    \item Sterilise bottles and dummies during infant oral thrush treatment
    \item Do not keep self-treating monthly without medical review---ask for examination and swabs
\end{itemize}

\section*{Sources and Further Reading}
\begin{itemize}
    \item NHS. \textit{Thrush in men and women.} https://www.nhs.uk/
    \item Mayo Clinic. \textit{Yeast infection (candidiasis).} https://www.mayoclinic.org/
    \item MSD Manual. \textit{Candidiasis.} https://www.msdmanuals.com/
    \item MedlinePlus. \textit{Yeast infections.} https://medlineplus.gov/
    \item Better Health Channel. \textit{Thrush.} https://www.betterhealth.vic.gov.au/
    \item Healthdirect. \textit{Thrush.} https://www.healthdirect.gov.au/
\end{itemize}
